\section{The generic state-space setting} \label{sec:setting}

This section pins down the abstractions \texttt{smc2fc} expects from the
user's model. Everything in subsequent sections depends only on these
two interfaces.

\subsection{The hidden Markov model in continuous time}

The framework assumes a continuous-time stochastic system observed at
discrete times. The latent state $X(t) \in \R^{n_x}$ evolves under a
controlled It\^o stochastic differential equation
\begin{equation}
\dd X(t) \;=\; b(X(t),\, U(t),\, \thetadyn)\, \dd t
            \;+\; \sigma(X(t),\, \thetadyn)\, \dd W(t),
\label{eq:sde}
\end{equation}
with drift $b: \R^{n_x} \times \R^{n_u} \times \R^d \to \R^{n_x}$,
diffusion $\sigma: \R^{n_x} \times \R^d \to \R^{n_x \times n_x}$
(typically diagonal in framework usage), Brownian motion $W(t)$, an
exogenous control input $U(t) \in \R^{n_u}$, and a static parameter
vector $\thetadyn \in \R^d$. Time is discretised on a uniform grid
$t_n = n \Delta t$, $n = 0, 1, 2, \ldots$, and equation~\eqref{eq:sde}
is integrated step-by-step by Euler--Maruyama:
\begin{equation}
X_{n+1} \;=\; X_n \;+\; b(X_n, U_n, \thetadyn) \Delta t
              \;+\; \sigma(X_n, \thetadyn) \sqrt{\Delta t}\, \xi_{n+1},
              \qquad \xi_{n+1} \sim \Ncal(0, I).
\label{eq:em}
\end{equation}
At each time step $n \ge 1$ the user obtains an observation
$Y_n \in \R^{n_y}$ via the user-supplied likelihood
\begin{equation}
Y_n \mid X_n \;\sim\; p\bigl(\,\cdot \mid X_n,\, C_n,\, \thetadyn\bigr),
\label{eq:obs}
\end{equation}
where $C_n$ is any per-step exogenous regressor (e.g.\ a circadian
forcing, a calendar feature) that the user wants to expose to the
observation model. The framework supports Gaussian, Bernoulli, and
mixed channels in $p(Y_n \mid X_n, \cdot)$ — all log-likelihoods are
summed across channels and any per-bin presence mask supplied by the
user.

The framework does not place semantic constraints on $X_n$, $U_n$, or
$Y_n$. It only requires that the drift, diffusion, observation
likelihood, and prior $p(\thetadyn)$ are all evaluable and
JAX-traceable.

\subsection{The two abstractions the framework consumes} \label{sec:interfaces}

\paragraph{(A) The forward simulator: \texttt{SDEModel}.}
A frozen dataclass that bundles the drift, diffusion, observation
samplers, and per-state physical bounds. Used by the user's
\texttt{StepwisePlant} (Section~\ref{sec:mpc}) for closed-loop
ground-truth simulation, and by the test infrastructure for synthetic
data generation. See \path{smc2fc/simulator/sde_model.py}.

\paragraph{(B) The inference target: \texttt{EstimationModel}.}
A frozen dataclass that bundles the propagate function, the diffusion
function, the observation log-weight function, and configuration of
the parameter prior. The framework's outer SMC$^2$ engine reads only
from this. See \path{smc2fc/estimation_model.py}.

\bigskip
The user supplies one of each. The two are typically thin wrappers
around the same underlying drift and observation code; the split exists
because the simulator side wants numpy-friendly per-step samplers
suitable for a Python loop, while the inference side wants pure-JAX
callables suitable for batching over many particles inside
\texttt{jax.vmap}.

\subsection{What the framework does \emph{not} require}

\begin{itemize}
  \item \emph{No closed-form transition density.} The framework only
        needs to be able to \emph{sample} from
        $p(X_n \mid X_{n-1}, U_n, \thetadyn)$, never to evaluate it
        pointwise. This is what makes it applicable to nonlinear SDEs.
  \item \emph{No tractable marginal likelihood.}
        $p(Y_{1:T} \mid \thetadyn)$ never has to be evaluated exactly;
        the bootstrap particle filter (Section~\ref{sec:pf}) provides
        an unbiased estimator that the outer SMC$^2$ then operates on.
  \item \emph{No Gaussianity.} Neither the transition kernel nor the
        observation likelihood needs to be Gaussian. The HMC moves in
        the outer SMC$^2$ act on the parameter vector $\thetadyn$, not
        on the latent state.
  \item \emph{No separation principle.} The framework does not assume
        the controller can be designed assuming a perfect state
        estimate. The receding-horizon loop of Section~\ref{sec:mpc}
        explicitly threads an SMC$^2$ posterior through the controller
        without separating estimation from control.
\end{itemize}

\subsection{Notational summary}

\begin{center}\small
\begin{tabular}{lll}
\toprule
Symbol & Meaning & Code \\
\midrule
$X_n \in \R^{n_x}$    & Latent state at step $n$                    & \texttt{state} (6D in FSA, 35D in SWAT) \\
$Y_n \in \R^{n_y}$    & Observation at step $n$                    & per-channel \texttt{obs\_value} arrays \\
$U_n,\, \Phi_n$       & Exogenous control input at step $n$        & \texttt{Phi} (sometimes \texttt{U}) \\
$C_n$                 & Per-step exogenous regressor               & \texttt{C} \\
$\thetadyn \in \R^d$  & Static parameter vector (inferred)         & \texttt{u} or \texttt{theta\_dyn} \\
$\thetactrl$          & Control-schedule parameter vector          & \texttt{theta\_ctrl} \\
$N$                   & Number of inner-PF state particles         & \texttt{n\_pf\_particles} \\
$M$                   & Number of outer-SMC parameter particles    & \texttt{n\_smc\_particles} \\
$\lambda$             & Tempering parameter, $0 \to 1$             & \texttt{lam} or \texttt{tempering\_param} \\
$\Lhat_N(\thetadyn)$  & Unbiased likelihood estimator              & \texttt{log\_density(u, ...)} \\
\bottomrule
\end{tabular}
\end{center}
